\documentclass[11pt]{article}
% ===== Shared preamble for Calendar notes =====
\usepackage[margin=1in]{geometry}
\usepackage{amsmath,amssymb}
\usepackage{enumitem}
\usepackage{xcolor}
\usepackage{tcolorbox}
\tcbuselibrary{skins,breakable}
\usepackage{titlesec}
\usepackage{fancyhdr}
\usepackage{hyperref}
\usepackage{array}
\usepackage{booktabs}
\usepackage{longtable}
\usepackage{pifont} % for \ding

% ---- Colors ----
\definecolor{primary}{HTML}{1B3A6B}
\definecolor{accent}{HTML}{2E7D32}
\definecolor{warnclr}{HTML}{B45309}
\definecolor{tipclr}{HTML}{0F5D8C}
\definecolor{lightbg}{HTML}{F4F6F8}

% ---- Section styling ----
\titleformat{\section}
  {\normalfont\Large\bfseries\color{primary}}
  {\thesection}{0.5em}{}[{\color{primary}\titlerule[1pt]}]
\titleformat{\subsection}
  {\normalfont\large\bfseries\color{primary}}
  {\thesubsection}{0.5em}{}

% ---- Header/Footer ----
\pagestyle{fancy}
\fancyhf{}
\renewcommand{\headrulewidth}{0.4pt}
\fancyfoot[C]{\thepage}

% ---- Hyperref ----
\hypersetup{
  colorlinks=true,
  linkcolor=primary,
  urlcolor=tipclr,
  citecolor=accent
}

% ---- Custom boxes ----
% Formula / Key concept box
\newtcolorbox{formulabox}[1][]{
  enhanced, breakable,
  colback=lightbg, colframe=primary,
  boxrule=0.8pt, arc=2pt,
  left=8pt, right=8pt, top=6pt, bottom=6pt,
  title={\textbf{#1}},
  fonttitle=\bfseries\color{white},
  colbacktitle=primary,
  attach boxed title to top left={xshift=8pt, yshift=-2mm},
  boxed title style={boxrule=0pt, arc=2pt}
}

% Tip / trick box
\newtcolorbox{tipbox}[1][Tip]{
  enhanced, breakable,
  colback=blue!4, colframe=tipclr,
  boxrule=0.8pt, arc=2pt,
  left=8pt, right=8pt, top=6pt, bottom=6pt,
  title={\textbf{\textcolor{white}{\ding{72}\ #1}}},
  fonttitle=\bfseries, colbacktitle=tipclr,
}

% Warning / common mistake box
\newtcolorbox{warnbox}[1][Watch Out]{
  enhanced, breakable,
  colback=orange!6, colframe=warnclr,
  boxrule=0.8pt, arc=2pt,
  left=8pt, right=8pt, top=6pt, bottom=6pt,
  title={\textbf{\textcolor{white}{#1}}},
  fonttitle=\bfseries, colbacktitle=warnclr,
}

% Example / worked question box
\newtcolorbox{examplebox}[1][Example]{
  enhanced, breakable,
  colback=green!4, colframe=accent,
  boxrule=0.8pt, arc=2pt,
  left=8pt, right=8pt, top=6pt, bottom=6pt,
  title={\textbf{\textcolor{white}{#1}}},
  fonttitle=\bfseries, colbacktitle=accent,
}


\title{\color{primary}\textbf{Calendar --- Fast Tricks \& Shortcuts}\\[4pt]\large Anchor-Day Method, Month Codes, and Every Speed Trick You Need}
\author{Aptitude Grind}
\date{}

\lhead{\small Calendar --- Fast Tricks \& Shortcuts}
\rhead{\small Aptitude Grind}

\begin{document}
\maketitle
\thispagestyle{fancy}

This file collects the \textbf{fast calculation methods} --- the actual shortcuts you execute under
exam pressure. For the full catalogue of question \emph{types} that use these tricks, see
\texttt{03-question-pattern-catalog.tex}.

\section{The Anchor-Day Method (Fastest for Placement Tests)}

Instead of counting odd days from year 1 (or even using the mod-400 shortcut from
\texttt{01-fundamentals.tex}), work relative to a \textbf{known reference date} close to your target
date. Most questions either give you a reference day directly, or you can use a memorized anchor.

\begin{formulabox}[Useful Anchors to Memorize]
\begin{itemize}[leftmargin=1.4em]
  \item 1 January 2000 $\to$ \textbf{Saturday} (most practical anchor for modern-date questions)
  \item 1 January 1900 $\to$ \textbf{Monday}
  \item 1 January 2100 $\to$ \textbf{Friday}
\end{itemize}
\end{formulabox}

\textbf{Anchor-Day Procedure}
\begin{enumerate}[leftmargin=1.6em]
  \item Find the number of complete years between the anchor and target year.
  \item Odd days contributed = (ordinary years $\times$ 1) + (leap years $\times$ 2), reduced mod 7.
  \item Add odd days for the months elapsed in the target year, then add the date number.
  \item Add this to the anchor's day-of-week position (mod 7) and read off the result.
\end{enumerate}

\begin{tipbox}[Why This Beats the Mod-400 Method]
The mod-400 method is for when you're asked about an arbitrary date, possibly centuries away. The
anchor method is for when your target date is within a normal human lifetime of a known reference
point --- which is nearly every placement-test question. Fewer years to sum = fewer chances to make
an arithmetic slip.
\end{tipbox}

\begin{examplebox}[Worked Example]
\emph{If 1 Jan 2000 was a Saturday, what day was 1 Jan 2010?}
\begin{itemize}[leftmargin=1.4em]
  \item Years 2000--2009 (10 years elapsed up to 1 Jan 2010)
  \item Leap years: 2000, 2004, 2008 $\to$ 3 leap, 7 ordinary
  \item Odd days $= (3\times2)+(7\times1) = 13 \equiv 6 \pmod 7$
  \item Saturday $+6 \to$ \textbf{Friday}
\end{itemize}
\end{examplebox}

\section{Month Codes Method (a.k.a. Doomsday-style Shortcut)}

A faster alternative used heavily in coaching materials: assign each month a fixed \textbf{code}
representing its odd-day contribution from 1 January of a non-leap year up to the 1st of that month.

\begin{formulabox}[Month Codes --- Non-Leap Year]
\centering
\begin{tabular}{lcccccccccccc}
\toprule
Month & Jan & Feb & Mar & Apr & May & Jun & Jul & Aug & Sep & Oct & Nov & Dec \\
\midrule
Code & 0 & 3 & 3 & 6 & 1 & 4 & 6 & 2 & 5 & 0 & 3 & 5 \\
\bottomrule
\end{tabular}
\end{formulabox}

\begin{tipbox}[Leap Year Adjustment]
For a \textbf{leap year}, subtract 1 from the codes of \textbf{January and February only} (the extra
day, 29 Feb, hasn't occurred yet when computing odd days up to those months). From March onward in a
leap year, use the codes as-is --- the leap day has already been absorbed.
\end{tipbox}

\[
\text{Day code} = (\text{Date} + \text{Month code} + \text{Year's odd days from anchor}) \bmod 7
\]
Match the result against the Odd-Days $\to$ Day table in \texttt{01-fundamentals.tex}.

\begin{warnbox}[When to Use Which Method]
The odd-days-from-scratch method is more intuitive and safer under exam pressure because you derive
everything instead of recalling a code table. The month-code method is faster once memorized, but a
single misremembered code silently gives a wrong answer. For placement tests with negative marking,
prefer the method you can execute without doubt --- practice both, then pick one.
\end{warnbox}

\section{Day After / Before N Days --- Pure Modulo, No Calendar Needed}

This is the single fastest calculation in the entire topic, and it needs \textbf{no leap-year
knowledge at all} when you're just counting raw days forward or backward from a \emph{named} day
(not a date).

\begin{formulabox}[Rule]
$N \bmod 7$ tells you how many weekday-steps to move. Forward for ``after,'' backward for ``before.''
\end{formulabox}

\begin{examplebox}[Example --- After]
Today is Tuesday. What day after 100 days? \\
$100 \bmod 7 = 2$ $\to$ Tuesday $+2 \to$ \textbf{Thursday}
\end{examplebox}

\begin{examplebox}[Example --- Before]
100 days before Friday. \\
$100 \bmod 7 = 2$ $\to$ Friday $-2 \to$ \textbf{Wednesday}
\end{examplebox}

\begin{warnbox}[Common Mistake]
Reaching for the odd-days-per-year method here. If the question only gives you a day name (not a
calendar date) and a day count, you never need leap years --- just $N \bmod 7$. Leap years only
matter once you're crossing actual calendar dates, because then the \emph{number of days in between}
depends on whether 29 February falls inside that span.
\end{warnbox}

\section{Leap Year --- Quick Check and Quick Count}

\begin{formulabox}[Quick Check for One Year]
Divisible by 4? No $\to$ Not leap. Yes $\to$ century year (ends in 00)? No $\to$ Leap. Yes $\to$
divisible by 400? Yes $\to$ Leap; No $\to$ Not leap.
\end{formulabox}

\begin{formulabox}[Quick Count for a Range]
Leap years from 1 to $N$ $= \lfloor N/4 \rfloor - \lfloor N/100 \rfloor + \lfloor N/400 \rfloor$. \\
For a range $[A,B]$: $\text{count}(B) - \text{count}(A-1)$.
\end{formulabox}

\begin{examplebox}[Worked Example --- Leap Years From 1601 to 2000]
\begin{align*}
\text{count}(2000) &= \lfloor 2000/4\rfloor - \lfloor 2000/100\rfloor + \lfloor 2000/400\rfloor = 500-20+5 = 485 \\
\text{count}(1600) &= \lfloor 1600/4\rfloor - \lfloor 1600/100\rfloor + \lfloor 1600/400\rfloor = 400-16+4 = 388 \\
\text{Leap years in }[1601,2000] &= 485-388 = \textbf{97}
\end{align*}
(Verified independently with a direct year-by-year count.)
\end{examplebox}

\section{Weekday Frequency Inside a Month}

Every month has at least 4 of each weekday. Whether a weekday occurs a 5th time depends purely on
\textbf{how many days the month has beyond 4 full weeks (28 days)}, and \textbf{which weekday the
month starts on}.

\begin{formulabox}[Rule]
Let $L = (\text{days in month}) - 28$. The \textbf{first $L$ weekdays} starting from the month's
1st-day weekday occur \textbf{5 times}; every other weekday occurs exactly 4 times.
\end{formulabox}

\begin{center}
\begin{tabular}{lcl}
\toprule
\textbf{Month length} & \textbf{L} & \textbf{Weekdays occurring 5 times} \\
\midrule
28 days (Feb, non-leap) & 0 & None --- every weekday occurs exactly 4 times \\
29 days (Feb, leap) & 1 & Only the weekday the month starts on \\
30 days & 2 & The starting weekday and the next one \\
31 days & 3 & The starting weekday and the next two \\
\bottomrule
\end{tabular}
\end{center}

\begin{examplebox}[Worked Example]
A 31-day month starts on Tuesday. Which weekdays occur 5 times? \\
$L = 31-28 = 3 \to$ first 3 weekdays from Tuesday: \textbf{Tue, Wed, Thu} \\
(Verified with a full day-by-day count for a real 31-day month starting on a Monday --- Mon/Tue/Wed
each appeared 5 times, matching the rule.)
\end{examplebox}

\section{Month Start/End Day --- Jump Straight From One to the Other}

\begin{formulabox}[Rule]
To go from the weekday of the 1st to the weekday of any day $D$ of the same month, add
$(D-1) \bmod 7$ to the 1st's weekday.
\end{formulabox}

\begin{examplebox}[Worked Example --- Start to End]
1 June = Wednesday. What day is 30 June? \\
$(30-1) \bmod 7 = 29 \bmod 7 = 1 \to$ Wednesday $+1 \to$ \textbf{Thursday}
\end{examplebox}

\begin{examplebox}[Worked Example --- Start of Next Month]
1 March = Monday. Find 1 April. (March has 31 days.) \\
$31 \bmod 7 = 3 \to$ Monday $+3 \to$ \textbf{Thursday} \\
This chains naturally: keep adding (month length mod 7) to hop from one month's 1st to the next.
\end{examplebox}

\section{Same-Calendar-Year Gaps --- Corrected}

Coaching material widely repeats ``an ordinary year repeats after 6, 11, \textbf{17}, or 28 years.''
This was checked directly by scanning three thousand years of real calendars, and \textbf{17 never
actually occurs}. Here's what's actually true:

\begin{center}
\begin{tabular}{lll}
\toprule
\textbf{Starting year type} & \textbf{Typical gap} & \textbf{Rare gap (broken century boundary)} \\
\midrule
Ordinary & 6 or 11 years & 12 years \\
Leap & 28 years & 40 years \\
\bottomrule
\end{tabular}
\end{center}

\begin{warnbox}[Don't Memorize ``17'']
It's a commonly repeated number in Indian test-prep material that does not hold up when checked
against actual calendars. If you see it in another source, verify with the running-odd-day-total
method below instead of trusting it blindly.
\end{warnbox}

\begin{formulabox}[Running-Odd-Day-Total Method --- Always Correct]
\begin{enumerate}[leftmargin=1.6em]
  \item Starting from your given year, add 1 odd day for each ordinary year that passes and 2 for
    each leap year, keeping a running total mod 7.
  \item Stop at the first year where the running total hits 0 \textbf{and} the year's leap/ordinary
    type matches the starting year's type (same weekday for 1 Jan isn't enough if one year is leap
    and the other isn't, since their Februaries differ).
\end{enumerate}
\end{formulabox}

\section{Quick Reference --- Which Trick for Which Situation}

\begin{center}
\begin{tabular}{p{7.7cm}l}
\toprule
\textbf{Situation} & \textbf{Use} \\
\midrule
Target date centuries away from any known anchor & Mod-400 method (Notes 01) \\
Target date within $\sim$100--150 years of a known anchor & Anchor-Day method (\S1) \\
You've memorized the month-code table & Month-code method (\S2) \\
Only a day name + a day count given (no calendar date) & $N \bmod 7$ (\S3) \\
``How many leap years between X and Y'' & Leap-count formula (\S4) \\
``Which weekday(s) occur 5 times'' & Weekday-frequency rule (\S5) \\
Need another day's weekday within the same month & Start/end day rule (\S6) \\
``Same calendar as year X repeats in which year'' & Running-odd-day-total (\S7) --- don't guess 17 \\
\bottomrule
\end{tabular}
\end{center}

\end{document}
